%%%%%%%%%%%%%%%%%%%%%%%%%%%%%%%%%
%
%       EXERCÍCIO
%
%%%%%%%%%%%%%%%%%%%%%%%%%%%%%%%%%

\ifdefstring{\atividade}{prova}{%
    \ifdefstring{\modo}{objetivo}{%
        \renewcommand{\valorquestao}{\ValorQObj\ ponto}
    }{%
        \renewcommand{\valorquestao}{\ValorQDisc\ pontos}
    }%
}%

\begin{exercicioBanco}[\valorquestao]
Em uma turma, \(40\) estudantes responderam a uma pesquisa sobre duas plataformas de streaming, \(A\) e \(B\).
Sabe-se que:

\begin{enumerate}[\(\bullet\)]
    \begin{multicols}{2}
        \item \(18\) estudantes utilizam a plataforma \(A\);
        \item \(9\) estudantes utilizam ambas as plataformas;
        \item \(11\) estudantes não utilizam nenhuma das duas plataformas.
    \end{multicols}
\end{enumerate}

Sabendo que todos os \(40\) estudantes responderam à pesquisa, o número de estudantes que utilizam a plataforma \(B\) é:

% Define as alternativas
\newcommand{\alternativas}{%
    \begin{center}
        \begin{tabularx}{\textwidth}{XXXXX}
            (a) \(11\). &
            (b) \(14\). &
            (c) \(17\). &
            (d) \(20\). &
            (e) \(22\).
        \end{tabularx}
    \end{center}
}

% Define a resposta correta
\newcommand{\resposta}{D}

% Lógica condicional para exibição
\ifdefstring{\atividade}{lista}{%
    \alternativas
    \vspace{0.5em}
    
    \noindent\textbf{Resposta:} letra \textbf{\resposta}.
}{%
    \ifdefstring{\modo}{objetivo}{%
        \alternativas
    }{%
        % modo = discursiva → não mostra alternativas
    }
}
\end{exercicioBanco}

